\chapter{Skin Biopsy}
\section*{What Is This Test?} Skin biopsy removes a small portion of skin for microscopic examination to diagnose inflammatory, infectious, pigmented, and cancerous lesions.\section*{Alternative Names} Punch biopsy; shave biopsy; excisional biopsy.\section*{Principle} Histopathology examines skin architecture and cells; special stains, immunofluorescence, or molecular tests may be added.\section*{Why This Test Is Ordered} It evaluates suspicious moles, persistent rashes, tumors, blisters, or lesions with uncertain diagnosis.\section*{Primary vs Secondary Test} Biopsy is a secondary diagnostic test after clinical examination and dermoscopy.\section*{How This Test Is Performed} After local anesthetic, tissue is shaved, punched, or excised, preserved, and sent to pathology.\section*{What Preparation Is Needed} Report anticoagulants, bleeding disorders, allergies, and medicines; do not stop medication without advice.\section*{Understanding Results} The report may identify benign, inflammatory, infectious, precancerous, or malignant change and whether margins are involved.\section*{Follow-up Tests} Special stains, immunohistochemistry, re-excision, imaging, or dermatology/oncology referral may follow.\section*{References} \begin{enumerate}\item MedlinePlus. Skin Biopsy.\item American Academy of Dermatology. Skin biopsy information.\end{enumerate}\clearpage\section*{Understanding Results and Follow-up}
The laboratory report should identify the specimen, method, units, reference interval or decision limit, and whether the result is positive, negative, abnormal, or inconclusive. Reference intervals vary with age, sex, pregnancy, specimen type, assay, and laboratory; a result outside a range is not automatically a diagnosis, and a result within a range does not always exclude disease. Discuss unexpected results with the ordering clinician, who can decide whether repeat or confirmatory testing, treatment, or referral is appropriate. Do not change medicines or delay urgent care based on an isolated result.
\section*{Regulatory Status and Authorization Date}This chapter describes a test category, not a specific commercial product. FDA, CDSCO, EU IVDR, and MHRA approval or registration dates therefore cannot be assigned without the assay manufacturer, product name, instrument, and intended use. For laboratory-developed tests, an individual agency approval date may not exist. See the Preface for the interpretation of regulatory status.
\clearpage
